% ---------------------------------------------------------------------
%  附录 C 文本格式定义
% ---------------------------------------------------------------------
\chapter{文本格式定义}
\label{app:formats}

本附录用简化的文法描述四种格式。记号：\code{int} 为十进制整数，
\code{float} 为浮点数，\code{name} 为不含空白的标识符，
\code{NL} 表示换行（newline），\code{SP} 表示空格（space），
\code{*} 表示零次或多次，\code{[...]} 表示可选。

\section{标准轮廓}
\begin{lstlisting}[numbers=none,language=plaintext]
file        := section+
section     := header param* "data" ":" block*
header      := "# " title NL
param       := key ":" SP value NL
key         := "sign" | "type" | "ntime" | "ntau"
             | "ndim1" | "ndim2" | "tmax" | "beta"
             | "dt" | "dtau"
block       := "index" ":" int SP* NL row*
row         := int SP int SP float SP float NL
\end{lstlisting}
整数索引从 $1$ 开始；\code{ntime}、\code{ntau} 是点数而不是
间隔数；推迟与小于段落写出完整方阵，未存储的三角由厄米关系
重构（非方阵时为零）。

\section{赝粒子}
与标准轮廓相同，只把对象头换成
\begin{lstlisting}[numbers=none,language=plaintext]
header := "# Pseudo-Particle Contour-Ordered Green's Function:" NL
\end{lstlisting}
读取端据此拒绝普通文件。

\section{移动窗}
\begin{lstlisting}[numbers=none,language=plaintext]
file        := header row*
header      := "# KB_MOVING_V1" SP kind SP int SP int
               SP int SP int SP int NL
kind        := "F" | "Cf"
row_F       := int SP int SP int SP int SP float SP float NL
row_Cf      := int SP int SP int SP float SP float NL
\end{lstlisting}
\code{F} 的行依次是 \code{delti deltj a b re im}，\code{Cf} 的行
是 \code{delti a b re im}。数据按逻辑行顺序写出。

\section{稳态}
\begin{lstlisting}[numbers=none,language=plaintext]
file        := header row*
header_fn   := "# SteadyFunction" SP int SP int NL
header_g    := "# SteadyGreen" SP int SP int SP float SP float NL
row         := dom SP comp SP int SP int SP int SP float SP float NL
dom         := "time" | "freq"
comp        := "ret" | "les"
\end{lstlisting}
每个文件必须包含全部模式、两个域与两个分量的数据；行数由头部
参数唯一确定。

\section{参考数据}
\begin{lstlisting}[numbers=none,language=plaintext]
file     := "# ref 1" NL meta* "# columns component i j a b re im" NL row*
meta     := "# " key SP value NL
row      := name SP int SP int SP int SP int SP float SP float NL
\end{lstlisting}
元数据的键包括 \code{data_kind}、\code{generator}、
\code{case_set}、\code{case}、\code{case_id}，以及各类型特有的
网格参数。所有索引从 $0$ 开始。
